\section*{Hoofdstuk \ref{ch:b3:palaeomagnetism} --- Geomagnetisme en paleomagnetisme}

\begin{solution}{exo:b3:palaeomagnetism:1}
$\tan I = 2\tan\ang{30} = 1.1547$, dus $I = \ang{49.1}$. Omgekeerd, $\tan\lambda
= \tfrac12\tan(\ang{-20}) = -0.1820$, dus $\lambda = \ang{-10.3}$, dat is
\ang{10.3}\,Z.
\end{solution}

\begin{solution}{exo:b3:palaeomagnetism:2}
$H = \sqrt{20000^2 + 2000^2} = \qty{20100}{nT}$;
$F = \sqrt{20000^2 + 2000^2 + 40000^2} = \qty{44770}{nT}$. De magnetische
declinatie is $\arctan(-2000/20000) = \ang{-5.7}$, \ang{5.7} naar het westen van het noorden;
de inclinatie $\arctan(40000/20100) = \ang{63.3}$, naar beneden.
\end{solution}

\begin{solution}{exo:b3:palaeomagnetism:3}
$B_0 = -g_1^0 = \qty{29350}{nT}$ aan de evenaar, en $2B_0 = \qty{58700}{nT}$ bij
de polen. $m = 4\pi a^3B_0/\mu_0 = a^3B_0/\num{e-7} = \num{2.5862e20}\times
\num{2.935e-5}/\num{e-7} = \qty{7.59e22}{A.m^2}$, iets minder dan de
\qty{7.69e22}{A.m^2} van de volledige dipool, die $g_1^1$ en $h_1^1$ toevoegt.
\end{solution}

\begin{solution}{exo:b3:palaeomagnetism:4}
(a) Een \omterm{def:b3:palaeomagnetism:nrm}{thermoremanente magnetisatie}, die het veld vastlegt toen de lava afkoelde
door de blokkeringstemperaturen van de korrels heen. (b) Een \omterm{def:b3:palaeomagnetism:nrm}{detritische remanente
magnetisatie}, die het veld vastlegt binnen de tijd die het silt nodig had om te bezinken en
door compactie vergrendeld te raken. (c) Een \omterm{def:b3:palaeomagnetism:nrm}{chemische remanente magnetisatie}, die het veld
vastlegt tijdens de groei van het hematiet, tijdens de vroege begraving en dus enigszins
na de afzetting. (d) Een viskeuze remanente magnetisatie langs het huidige veld, toegevoegd aan de
oorspronkelijke TRM door de korrels met de kortste relaxatietijd; zij wordt verwijderd door de
eerste stappen van de demagnetisatie.
\end{solution}

\begin{solution}{exo:b3:palaeomagnetism:5}
De eenheidsvectoren $(\cos I\cos D,\ \cos I\sin D,\ \sin I)$ zijn
\begin{gather*}
  (0.6545, 0.1391, 0.7431),\quad (0.6670, 0.2695, 0.6947),\\
  (0.5235, 0.1501, 0.8387),\quad (0.5985, 0.0629, 0.7986).
\end{gather*}
De som is $(2.4435, 0.6216, 3.0751)$, met lengte $R = 3.9766$. Gemiddelde:
$D = \arctan(0.6216/2.4435) = \ang{14.3}$, $I = \arcsin(3.0751/3.9766) =
\ang{50.7}$. $k = 3/(4 - 3.9766) = 128$. Met $20^{1/3} - 1 = 1.7144$,
$\cos\alpha_{95} = 1 - (0.0234/3.9766)\times1.7144 = 0.98991$, dus
$\alpha_{95} = \ang{8.1}$: met vier monsters blijft de kegel breed, hoe
strak de groep ook is.
\end{solution}

\begin{solution}{exo:b3:palaeomagnetism:6}
\sloppy $\lambda = \arctan(\tfrac12\tan\ang{50}) = \ang{30.79}$, dus $p = \ang{59.21}$.
$\sin\lambda_p = \sin\ang{40}\cos\ang{59.21} + \cos\ang{40}\sin\ang{59.21}
\cos\ang{20} = 0.9474$, $\lambda_p = \ang{71.34}$\,NB. $\sin\beta =
\sin\ang{59.21}\sin\ang{20}/\cos\ang{71.34} = 0.9184$, $\beta = \ang{66.66}$.
Omdat $\cos p = 0.5119 < \sin\ang{40}\times0.9474 = 0.6090$, is het tweede geval
van toepassing: $\varphi_p = 10 + 180 - 66.66 = \ang{123.34}$\,OL. De pool ligt op
\ang{71.3}\,NB, \ang{123.3}\,OL.\par
\end{solution}

\begin{solution}{exo:b3:palaeomagnetism:7}
(a) $(7.69 - 8.32)/8.32 = \qty{-7.6}{\percent}$ in 125 jaar, of
\qty{-6.1}{\percent} per eeuw. (b) $29350/12.6 = 2329$, ongeveer 2300 jaar.
Het is geen voorspelling: de snelheid zelf verandert van het ene vijfjaarmodel naar het
volgende; het moment is in het verleden zowel hoger als lager geweest zonder
om te keren, en het veld keert op onregelmatige tijdstippen om die
geen enkele snelheid voorspelt.
\end{solution}

\begin{solution}{exo:b3:palaeomagnetism:8}
Tussen 20 en \qty{300}{\celsius} verloor de vector $(6.4, 0.0, 13.6)$: een
component van \qty{15.0}{mA/m}, met $D = \ang{0.0}$ en
$I = \arctan(13.6/6.4) = \ang{64.8}$. Boven \qty{300}{\celsius} verloor hij $(25.8, 18.1, 24.6)$:
\qty{40.0}{mA/m}, $D = \arctan(18.1/25.8) = \ang{35.1}$, $I =
\arctan(24.6/31.5) = \ang{38.0}$. De tweede, die naar de oorsprong gaat, is de
ChRM; de eerste, langs een steile, noordwaarts gerichte richting, bij lage
temperatuur verwijderd, is een viskeuze overprint langs het recente veld.
\end{solution}

\begin{solution}{exo:b3:palaeomagnetism:9}
$\lambda = \arctan(\tfrac12\tan(\ang{-35})) = \ang{-19.3}$. Voor $I = \ang{-40}$
en \ang{-30}: \ang{-22.8} en \ang{-16.1}, een interval van \ang{6.7}.
$p = \ang{90} - \lambda = \ang{109.3}$, $\cos^2p = 0.1096$:
$\mathrm dp = 2.5\times(1 + 0.329) = \ang{3.3}$ en $\mathrm dm =
5\times\sin\ang{109.3}/\cos\ang{35} = \ang{5.8}$.
\end{solution}

\begin{solution}{exo:b3:palaeomagnetism:10}
(a) $\tan\lambda = \tfrac12$, $\lambda = \ang{26.6}$. $F$ is $B_0$ aan de
evenaar en $2B_0$ bij de polen: een verhouding van 2. (b) $\tan I = 2\tan\lambda$
levert $(1 + \tan^2I)\,\mathrm dI = 2(1 + \tan^2\lambda)\,\mathrm d\lambda$, dus
\[
  \frac{\mathrm dI}{\mathrm d\lambda} = \frac{2/\cos^2\lambda}{1 + 4\tan^2\lambda}
  = \frac{2}{\cos^2\lambda + 4\sin^2\lambda} = \frac{2}{1 + 3\sin^2\lambda}.
\]
Een fout $\delta I$ levert $\delta\lambda = \tfrac12(1 + 3\sin^2\lambda)\,\delta I$,
van $\tfrac12\delta I$ aan de evenaar tot $2\,\delta I$ bij de polen: de
grootste fout ligt nabij de polen. (c) Nabij de evenaar verandert de inclinatie
twee keer zo snel als de breedtegraad, dus een gegeven fout in de inclinatie kost
de helft zoveel breedtegraad; nabij de polen verandert de inclinatie nauwelijks met de breedtegraad.
\end{solution}

\begin{solution}{exo:b3:palaeomagnetism:11}
(a) $\cos p = \sin(\ang{-35})\sin\ang{70} + \cos\ang{35}\cos\ang{70}\cos\ang{150}
= -0.5390 - 0.2426 = -0.7816$, $p = \ang{141.41}$. Dan $\tan I_x = 2\cot p =
-2.506$, $I_x = \ang{-68.2}$, en $\lambda_x = \ang{90} - p = \ang{-51.4}$.
$\cos D_x = (\sin\ang{70} - \sin(\ang{-35})\cos p)/(\cos\ang{35}\sin p) =
0.9617$, en omdat de pool oostelijk ligt van de meridiaan van de locatie ($\sin\ang{150} >
0$), $D_x = \ang{15.9}$. (b) $\lambda_o = \arctan(\tfrac12\tan(\ang{-60})) =
\ang{-40.9}$. $R = 30.0 - 15.9 = \ang{14.1}$: een rotatie van het blok met de klok mee
rond een verticale as. $P = \lambda_x - \lambda_o = -51.4 + 40.9 =
\ang{-10.5}$: de formatie lag \ang{10.5} verder naar het noorden dan de stabiele
plaat aangeeft dat ze zou moeten, dus het blok is sindsdien ongeveer \ang{10.5} verschoven, zo'n
\qty{1170}{km} naar het zuiden ten opzichte van de plaat.
\end{solution}

\begin{solution}{exo:b3:palaeomagnetism:12}
(a) Met $\mathrm d\Omega = 2\pi\sin\theta'\,\mathrm d\theta'$,
\[
  \int_0^\theta\frac{\kappa e^{\kappa\cos\theta'}}{4\pi\sinh\kappa}\,2\pi\sin\theta'\,
  \mathrm d\theta' = \frac{e^\kappa - e^{\kappa\cos\theta}}{2\sinh\kappa}
  = \frac{e^\kappa - e^{\kappa\cos\theta}}{e^\kappa - e^{-\kappa}} .
\]
(b) Voor grote $\kappa$ is $e^{-\kappa}$ verwaarloosbaar; delen door $e^\kappa$
levert $1 - e^{-\kappa(1 - \cos\theta)} \approx 1 - e^{-\kappa\theta^2/2}$. Dit
is gelijk aan $1 - e^{-1} = \qty{63}{\percent}$ wanneer $\kappa\theta^2/2 = 1$, $\theta =
\sqrt{2/\kappa}$ radialen, oftewel $\ang{81}/\sqrt\kappa$. (c) $81/\sqrt{50} =
\ang{11.5}$; de exacte formule geeft \qty{63.1}{\percent}.
\end{solution}

\begin{solution}{pb:b3:palaeomagnetism:1}
\textbf{1.} Hematiet in een rode laag kan groeien uit ijzeroxiden en hydroxiden
lang nadat het sediment is afgezet, waardoor een \omterm{def:b3:palaeomagnetism:nrm}{chemische remanente magnetisatie} van die
latere ouderdom ontstaat. Een positieve plooitest, als de lagen kort na
de afzetting zijn geplooid, of een positieve omkeringstest over dicht bij elkaar
liggende lagen van beide polariteiten zou aantonen dat de remanente magnetisatie vroeg is.

\textbf{2.} $(\cos\ang{37}\cos\ang{50}, \cos\ang{37}\sin\ang{50}, \sin\ang{37})
= (0.5134, 0.6118, 0.6018)$.

\textbf{3.} $R = \sqrt{3.8943^2 + 2.6500^2 + 3.5717^2} = 5.9114$. $D =
\arctan(2.6500/3.8943) = \ang{34.2}$, $I = \arcsin(3.5717/5.9114) = \ang{37.2}$.

\textbf{4.} $k = 5/(6 - 5.9114) = 56.5$. Met $20^{1/5} - 1 = 0.8206$,
$\cos\alpha_{95} = 1 - (0.0886/5.9114)\times0.8206 = 0.98771$, $\alpha_{95} =
\ang{9.0}$.

\textbf{5.} B: $k = 6/(7 - 6.8979) = 58.8$; $20^{1/6} - 1 = 0.6475$,
$\cos\alpha_{95} = 1 - (0.1021/6.8979)\times0.6475 = 0.99042$, $\alpha_{95} =
\ang{7.9}$. C: $k = 7/(8 - 7.8707) = 54.1$; $20^{1/7} - 1 = 0.5340$,
$\cos\alpha_{95} = 1 - (0.1293/7.8707)\times0.5340 = 0.99123$, $\alpha_{95} =
\ang{7.6}$.

\textbf{6.} B heeft de grootste $k$, de strakste groepering; C heeft de kleinste
$\alpha_{95}$, het best bepaalde gemiddelde, omdat het de meeste vindplaatsen telt. De drie
liggen dicht bij elkaar: alle drie formaties zijn goed gegroepeerd.

\textbf{7.} $\lambda_A = \arctan(\tfrac12\tan\ang{37.17}) = \ang{20.76}$\,N,
$\lambda_B = \ang{32.15}$\,N, $\lambda_C = \ang{45.33}$\,N.

\textbf{8.} $p = 90 - 20.76 = \ang{69.24}$. $\sin\lambda_p =
\sin\ang{47}\cos\ang{69.24} + \cos\ang{47}\sin\ang{69.24}\cos\ang{34.23} =
0.2593 + 0.5272 = 0.7865$: $\lambda_p = \ang{51.86}$\,N.

\textbf{9.} $\sin\beta = \sin\ang{69.24}\sin\ang{34.23}/\cos\ang{51.86} =
0.8516$, $\beta = \ang{58.40}$. Omdat $\cos p = 0.3545 < \sin\ang{47}\times
0.7865 = 0.5752$, geldt het tweede geval, $\cos\beta < 0$:
$\varphi_p = 15.00 + 180 - 58.40 = \ang{136.60}$\,O.

\textbf{10.} $\cos^2p = 0.1257$: $\mathrm dp = 4.50\times1.377 = \ang{6.2}$;
$\mathrm dm = 8.99\times\sin\ang{69.24}/\cos\ang{37.17} = \ang{10.5}$.

\textbf{11.} Van \ang{51.9}\,N bij \qty{60}{Ma} naar \ang{66.5}\,N bij
\qty{30}{Ma} en \ang{87.1}\,N bij \qty{4}{Ma}, op bijna dezelfde meridiaan
(\ang{136}--\ang{139}\,O): een pad richting de geografische pool, binnen zijn
kegel bereikt door de jongste formatie.

\textbf{12.} De aslaag, op \qty{22}{m}, ligt in de tweede normale zone,
19.7--\qty{24.6}{m}; met een ouderdom nabij \qty{3.2}{Ma} is die zone C2An.2n
(3.207--3.116). Dan is de zone 4.2--\qty{14.8}{m} C2An.3n, de zone
29.2--\qty{54.0}{m} C2An.1n, de omgekeerde basis hoort bij het chron onder
C2An, en de omgekeerde top hoort bij het chron erboven.

\textbf{13.} $(19.7 - 4.2)/(3.596 - 3.207) = 15.5/0.389 = \qty{39.8}{m/Myr}$;
$(54.0 - 24.6)/(3.116 - 2.581) = 29.4/0.535 = \qty{55.0}{m/Myr}$.

\textbf{14.} Basis: $3.596 + 4.2/39.8 = \qty{3.70}{Ma}$. Top: $2.581 -
10.0/55.0 = \qty{2.40}{Ma}$.

\textbf{15.} Alleen omkeringsgrenzen worden gedateerd, een paar honderdduizend jaar
uit elkaar; ertussen wordt de snelheid constant verondersteld. Een verandering van de subsidentie, van
de sedimenttoevoer of van het klimaat, of een hiaat in de afzetting, zou haar veranderen, en een
hiaat zou zelfs een korte zone kunnen wegnemen en de match misleiden.

\textbf{16.} $\cos p = \sin\ang{47}\sin\ang{78} + \cos\ang{47}\cos\ang{78}
\cos\ang{175} = 0.7154 - 0.1413 = 0.5741$, $p = \ang{54.96}$. $\tan I_x =
2\cot p = 1.4023$, $I_x = \ang{54.5}$; $\lambda_x = \ang{35.04}$\,N.

\textbf{17.} $\cos D_x = (\sin\ang{78} - \sin\ang{47}\times0.5741)/
(\cos\ang{47}\sin\ang{54.96}) = 0.99976$; de pool ligt oost van de
meridiaan van de vindplaats ($\sin\ang{175} > 0$), dus $D_x = \ang{1.3}$.

\textbf{18.} Verwachte \omterm{def:b3:palaeomagnetism:vgp}{paleobreedtes} uit $I_x$ = \ang{54.5}, \ang{60.0} en
\ang{63.6}: \ang{35.03}, \ang{40.89} en \ang{45.21}. $P_A = 35.03 - 20.76 =
\ang{14.3}$, $P_B = 40.89 - 32.15 = \ang{8.7}$, $P_C = 45.21 - 45.33 =
\ang{-0.1}$.

\textbf{19.} $R_A = 34.2 - 1.3 = \ang{32.9}$, $R_B = 23.9 - 2.9 = \ang{21.0}$,
$R_C = 3.4 - 2.5 = \ang{0.9}$, alle met de klok mee.

\textbf{20.} $\Delta D_o = \arcsin(\sin\ang{8.99}/\cos\ang{37.17}) = \ang{11.3}$;
$\Delta D_x = \arcsin(\sin\ang{3.0}/\cos\ang{54.5}) = \ang{5.2}$; $\Delta R_A =
\sqrt{11.3^2 + 5.2^2} = \ang{12.4}$.

\textbf{21.} $R_A = \ang{32.9} \pm \ang{12.4}$ en $R_B = \ang{21.0} \pm
\ang{13.8}$ zijn significant; $R_C = \ang{0.9} \pm \ang{17.9}$ is niet.

\textbf{22.} Zestig miljoen jaar geleden lag Kesta ongeveer \ang{14} breedtegraad
ten zuiden van waar het stabiele continent het plaatst, en was het ongeveer \ang{33}
tegen de klok in gedraaid ten opzichte van zijn huidige oriëntatie. Tegen \qty{30}{Ma} was het
\ang{6} van die \ang{14} naar het noorden opgeschoven en \ang{12} van die \ang{33}
met de klok mee gedraaid; tegen \qty{4}{Ma} had het zijn huidige plaats op het continent bereikt,
zonder rotatie over. Het dreef langs de marge van het continent naar het noorden en draaide
daarbij met de klok mee, en was tegen \qty{4}{Ma} op zijn plaats.

\textbf{23.} $47.00 - 20.76 = \ang{26.24}$, oftewel $26.24\times111.19 =
\qty{2918}{km}$. Dit omvat de beweging van het stabiele continent zelf;
ten opzichte daarvan heeft Kesta zich $P_A = \ang{14.3}$ verplaatst.

\textbf{24.} $I = \ang{28.18}$ en \ang{46.16} geven $\lambda = \ang{15.00}$ en
\ang{27.50}: een halve breedte van \ang{6.25}, en één standaarddeviatie van
$6.25/1.96 = \ang{3.19}$, oftewel \qty{355}{km}. Sinds formatie A is Kesta
naar het noorden gedreven met
\textbf{\qty[text-series-to-math=true]{2918 \pm 355}{km}}.
\end{solution}
